\section{相关工作}

本文位于思想史、政治思想、历史网络分析与大规模文本研究四条文献脉络的交叉处。它不以图谱替代思想史，而是追问一种通俗叙事把思想写成哪些人物关系与政治节点。

\subsection{良知、德性与行动}

王阳明研究首先处理“良知”作为道德知识的性质。Huang 将良知解释为一种新儒家智慧概念，强调道德判断不是外加规则，而是主体可发动的知识\cite{huang2006}；Tien 则从道德形上学与信念正当化两端讨论王阳明，指出情感、认识与道德实践不能被拆成互不相干的层次\cite{tien2004,tien2010}。这些研究面对的是哲学论证本身，材料来自思想文本及其概念脉络。本文面对的材料不同：通俗叙事把“良知”压缩成人物语言、师承与行动场景。因此，我们不判定其阐释是否穷尽原义，只测量概念在叙事中的出现及其所连接的社会关系。

\subsection{知行关系与政治秩序}

把心学理解为行动伦理，是第二条脉络。Angle 以德性伦理学讨论王阳明，把修养、判断与行动放进同一人格结构\cite{angle2010}；Yao 通过与 Bacon 的比较，强调学习并非知识囤积，而须在行动中完成\cite{yao2013}。与此同时，Kim 对王阳明与湛若水的争论表明，新儒家并不只讨论内在修养，也触及政治统一与秩序构造\cite{kim2012}。这组文献提醒我们，不能把“思想”预设为朝廷之外的纯观念空间。本文因此把王守仁的师生同门边与君臣、攻讦、冲突、同盟边放在同一分母内观察，而不是先把他归入“思想家”后便删去政治关系。

\subsection{中心性与历史关系网络}

网络分析提供了第三条方法线索。Freeman 对中心性概念的澄清表明，度数只刻画直接连接数量，不能自动等同于权力、声望或因果影响\cite{freeman1978}。Lemercier 进一步说明，历史网络方法的价值在于把关系选择、时间边界与资料缺失公开，而非让可视化图形替历史学家代言\cite{lemercier2015}。本文遵守这一限制：度数50只说明王守仁在2324条抽取关系中的直接连接数，不等同于“思想影响力50”；类别占比说明边如何被编码，也不证明行为动机。

\subsection{远读与可回到原文的核验}

大规模语料研究展示了从文本集合中提取趋势的可能性\cite{michel2011}，但单部著作的高频词尤其容易被误读。词频会把书名、章标题、引语与作者评论一并计数；同一名字的反复出现也可能来自情节密度，而非历史重要性。本文采用“两步远读”：先用逐字检索确认“心学”“知行合一”“讲学”“书院”等词的非零命中，再回到结构化事件和关系边。文本命中是存在性证据，图谱统计是结构证据，两者不互相冒充。

\subsection{本文的位置}

既有哲学研究解释王阳明思想“是什么”，网络研究解释关系量“如何算”，远读研究解释文本模式“怎样抽取”。本文把三者收束到一个较窄的问题：一部通俗明史怎样把思想传播写成师承，又怎样把同一批人物写入朝廷权力关系。研究对象不是完整明代思想世界，而是该书可观察、可复算的叙事结构；语料版本记为无 DOI 的书籍材料\cite{dangnian2009}。
